% formulas.tex -- Module 8 ESSENTIALS: All formulas at a glance.
% One row per formula of the module, in chapter order. Formulas, notation and sheet status
% are those of the full book's src/mod8/formulasheet.tex, which follows
%   mass defect Delta m = (Z m_p + N m_n) - m_nucleus: NOT ON SHEET;
%   binding energy E_B = Delta m c^2, E (MeV) = Delta m (u) x 931.5: COMBINE (as on the
%   Essentials ch1 and ch4 cards).
% One extra row: the photon momentum p = h/lambda, stated "not on the formula sheet" in the
% full book's ch3 and in the Essentials ch3.
% "Ch." = the chapter of this book where the formula is used.
% Rows with the same sheet status that belong together are joined into one row (as on the
% Essentials cards): v = E/B with q/m = E/(B^2 r); q = mgd/V with q = ne; mvr = nh/(2 pi)
% with 2 pi r = n lambda; N_t = N_0 e^(-lambda t) with lambda = ln 2 / t_half;
% r = mv/(qB), T = 2 pi m/(qB) with f = qB/(2 pi m).
% Data-sheet values and notes: copied from the full book's formulasheet.tex (values exactly
\unnumberedchapter[\faIcon{square-root-alt}]{All formulas at a glance}
\phantomsection\label{fg:chap}

% ---- helpers for this part only (prefix fg)
\newcommand\fgYes{\textcolor{hscformula}{\footnotesize\bfseries\faCheck~Yes}}
\newcommand\fgNo{\textcolor{hscmistake}{\footnotesize\bfseries\faTimes~Learn~it}}
\newcommand\fgComb{\textcolor{hsctiptext}{\footnotesize\bfseries\faLink~Combine}}
% a formula with 1.8pt of clear space above and below it, so fractions never touch the rules
\newcommand\fgf[1]{\raisebox{0pt}[\dimexpr\height+1.8pt\relax][\dimexpr\depth+1.8pt\relax]{#1}}
% \fgrow{mark}{formula}{what it is for}{chapter(s)}
\newcommand\fgrow[4]{\fgf{#2} & #3 & #1 & #4\\\hline}
% \fgrowii{mark}{formula, line 1}{formula, line 2}{what it is for}{chapter(s)}
\newcommand\fgrowii[5]{\fgf{#2}\newline\fgf{#3} & #4 & #1 & #5\\\hline}
% a darker rule between two chapters
\newcommand\fgnext{\noalign{\vskip-0.4pt{\color{hscbrand!45}\hrule height 0.9pt}}}
% chapter number, linked to the chapter
\newcommand\fgch[1]{\hyperref[ch#1:chap]{\textbf{#1}}}
% a group heading row
\newcommand\fggroup[1]{\multicolumn{4}{@{}l@{}}{\rule[-1.0ex]{0pt}{3.7ex}\normalsize
  {\bfseries\color{hscbrand}#1}}\\\hline}
\newcommand\fghead{\colhead{Formula} & \colhead{What it is for} & \colhead{Sheet} &
  \colhead{Ch.}}

Every formula in Module~\modulenumber, in chapter order, written as on the official formula
sheet. \fgYes: printed on the sheet, so you are given it in the exam. \fgNo: not on the
sheet, so learn it or learn how to derive it. \fgComb: not printed in this form; join
formulas from the sheet in one or two lines. ``Ch.'' is the chapter.

{\small
\setlength\LTpre{3pt}\setlength\LTpost{4pt}
\renewcommand\arraystretch{1.05}\arrayrulecolor{hscrule}
\begin{longtable}{@{}C{47mm}@{\hspace{3mm}}L{\dimexpr\linewidth-47mm-17.5mm-10mm-9mm\relax}%
    @{\hspace{3mm}}C{17.5mm}@{\hspace{3mm}}C{10mm}@{}}
  \fghead\\\hline\endfirsthead
  \fghead\\\hline\endhead
  % ---- Chapter 1: Origins of the Elements
  \fggroup{\Cref{ch1:chap}: Origins of the Elements}
  \fgrow{\fgYes}{$E = mc^2$}
    {Mass and energy; a pair needs $2mc^2$.}{\fgch{1},\,\fgch{4},\,\fgch{5}}
  \fgrow{\fgNo}{$z = \dfrac{\Delta\lambda}{\lambda_0} = \dfrac{\lambda - \lambda_0}{\lambda_0}$,
      \ $z \approx \dfrac{v}{c}$}
    {Redshift and speed of a galaxy.}{\fgch{1}}
  \fgrow{\fgNo}{$v = H_0 d$ \qquad $t \approx \dfrac{1}{H_0}$}
    {Hubble's law; rough age of the Universe.}{\fgch{1}}
  \fgrow{\fgYes}{$\lambda_\mathrm{max} = \dfrac{b}{T}$}
    {Black-body temperature (K) from its peak.}{\fgch{1}}
  \fgrow{\fgNo}{$\dfrac{L_A}{L_B} = 100^{(M_B - M_A)/5}$}
    {Compare luminosities from magnitudes.}{\fgch{1}}
  \fgrowii{\fgComb}{$E = \Delta m\,c^2$}{$E\,(\mathrm{MeV}) = \Delta m\,(\mathrm{u}) \times 931.5$}
    {Energy from the mass lost in a nuclear reaction; also the binding
     energy $E_\mathrm{B}$.}{\fgch{1}, \fgch{4}}
  \fgrow{\fgComb}{$\dfrac{\Delta m}{\Delta t} = \dfrac{P}{c^2}$}
    {Mass a star loses each second.}{\fgch{1}}
  \fgnext
  % ---- Chapter 2: Structure of the Atom
  \fggroup{\Cref{ch2:chap}: Structure of the Atom}
  \fgrow{\fgNo}{$v = \dfrac{E}{B}$ \qquad $\dfrac{q}{m} = \dfrac{E}{B^2 r}$}
    {Thomson: balanced fields, then $B$ alone.}{\fgch{2}}
  \fgrow{\fgNo}{$\dfrac{q}{m} = \dfrac{2V}{B^2 r^2}$}
    {School e/m tube (from rest through $V$).}{\fgch{2}}
  \fgrow{\fgNo}{$q = \dfrac{mgd}{V}$ \qquad $q = ne$}
    {Millikan's still drop; quantised charge.}{\fgch{2}}
  \fgrow{\fgYes}{$F = \dfrac{1}{4\pi\varepsilon_0}\dfrac{q_1 q_2}{r^2}$}
    {Force between an alpha and a nucleus.}{\fgch{2}}
  \fgrowii{\fgYes}{$\sum m\vec{v}_\mathrm{before} = \sum m\vec{v}_\mathrm{after}$}
    {$\sum\dfrac{1}{2}mv^2_\mathrm{before} = \sum\dfrac{1}{2}mv^2_\mathrm{after}$}
    {Chadwick's collisions (kinetic energy: elastic collisions only).}
    {\fgch{2}}
  \fgnext
  % ---- Chapter 3: The Quantum Atom
  \fggroup{\Cref{ch3:chap}: The Quantum Atom}
  \fgrow{\fgYes}{$E = hf$}
    {Energy of one photon.}{\fgch{3}}
  \fgrow{\fgComb}{$E = \dfrac{hc}{\lambda}$}
    {Photon energy from its wavelength.}{\fgch{3}}
  \fgrow{\fgComb}{$hf = E_\mathrm{i} - E_\mathrm{f}$}
    {Photon from a drop between two levels.}{\fgch{3}}
  \fgrow{\fgNo}{$E_n = -\dfrac{\qty{13.6}{eV}}{n^2}$}
    {Energy levels of hydrogen.}{\fgch{3}}
  \fgrow{\fgYes}{$\dfrac{1}{\lambda} = R\left(\dfrac{1}{n_\mathrm{f}^2} - \dfrac{1}{n_\mathrm{i}^2}\right)$}
    {Any hydrogen line (Balmer: $n_\mathrm{f} = 2$).}{\fgch{3}}
  \fgrow{\fgNo}{$mvr = \dfrac{nh}{2\pi}$ \qquad $2\pi r = n\lambda$}
    {Bohr's allowed orbits (whole waves fit).}{\fgch{3}}
  \fgrow{\fgNo}{$p = \dfrac{h}{\lambda}$}
    {Momentum of a photon.}{\fgch{3}}
  \fgrow{\fgYes}{$\lambda = \dfrac{h}{mv}$}
    {Matter waves (below about $0.1c$).}{\fgch{3}, \fgch{5}}
  \fgnext
  % ---- Chapter 4: Properties of the Nucleus
  \fggroup{\Cref{ch4:chap}: Properties of the Nucleus}
  \fgrowii{\fgNo}{$\sum A_\mathrm{before} = \sum A_\mathrm{after}$}
    {$\sum Z_\mathrm{before} = \sum Z_\mathrm{after}$}
    {Write or check a nuclear equation; find a missing nucleus.}{\fgch{4}}
  \fgrow{\fgYes}{$N_\mathrm{t} = N_0 e^{-\lambda t}$ \qquad $\lambda = \dfrac{\ln 2}{t_{\frac{1}{2}}}$}
    {Amount left after $t$; the decay constant.}{\fgch{4}}
  \fgrow{\fgComb}{$t = \dfrac{1}{\lambda}\ln\dfrac{N_0}{N_\mathrm{t}}$}
    {Time to fall from $N_0$ to $N_\mathrm{t}$.}{\fgch{4}}
  \fgrow{\fgNo}{$N_\mathrm{t} = N_0\left(\tfrac{1}{2}\right)^{n}$, \ $n = \dfrac{t}{t_{\frac{1}{2}}}$}
    {A whole number of half-lives.}{\fgch{4}}
  \fgrow{\fgNo}{$A = \lambda N$ \qquad $A_\mathrm{t} = A_0 e^{-\lambda t}$}
    {Activity in Bq ($\lambda$ in \si{s^{-1}}).}{\fgch{4}}
  \fgrow{\fgNo}{$\Delta m = (Zm_\mathrm{p} + Nm_\mathrm{n}) - m_\mathrm{nucleus}$}
    {Mass defect; $E_\mathrm{B} = \Delta m\,c^2$; per nucleon $E_\mathrm{B}/A$.}
    {\fgch{4}}
  \fgrow{\fgNo}{$E = \sum E_\mathrm{B,after} - \sum E_\mathrm{B,before}$}
    {Fission or fusion from binding energies.}{\fgch{4}}
  \fgnext
  % ---- Chapter 5: Deep Inside the Atom
  \fggroup{\Cref{ch5:chap}: Deep Inside the Atom}
  \fgrow{\fgYes}{$p_\mathrm{v} = \dfrac{m_0 v}{\sqrt{\left(1 - \dfrac{v^2}{c^2}\right)}}$}
    {Momentum faster than about $0.1c$.}{\fgch{5}}
  \fgrow{\fgNo}{$\lambda = \dfrac{h}{p_\mathrm{v}}$ \qquad $r = \dfrac{p_\mathrm{v}}{qB}$}
    {$\lambda$ and $r$ faster than about $0.1c$.}{\fgch{3}, \fgch{5}}
  \fgrow{\fgNo}{$Q = q_1 + q_2 + q_3$ \quad $Q = q_1 + q_2$}
    {Charge of a baryon or a meson.}{\fgch{5}}
  \fgrow{\fgYes}{$W = qV$}
    {Energy from one gap; $n$ gaps: $K = nqV$.}{\fgch{5}}
  \fgrowii{\fgNo}{$r = \dfrac{mv}{qB}$ \qquad $T = \dfrac{2\pi m}{qB}$}
    {$f = \dfrac{1}{T} = \dfrac{qB}{2\pi m}$}
    {A circle at \ang{90} to $B$ (well below $c$); the cyclotron's AC
     frequency.}{\fgch{5}}
  \fgnext
  % ---- From Modules 6 and 7, Year 11
  \fggroup{Other formulas used in this module (Modules 6 and 7, Year 11)}
  \fgrow{\fgYes}{$v = f\lambda$}
    {Wave equation; for light, $c = f\lambda$.}{\fgch{1}, \fgch{3}}
  \fgrow{\fgYes}{$d\sin\theta = m\lambda$}
    {A diffraction grating.}{\fgch{3}}
  \fgrow{\fgYes}{$P = \dfrac{\Delta E}{\Delta t}$}
    {Power of a star.}{\fgch{1}}
  \fgrow{\fgYes}{$E = \dfrac{V}{d}$ \qquad $\vec{F} = q\vec{E}$}
    {Field between plates; force on a charge.}{\fgch{2}}
  \fgrow{\fgYes}{$F = qv_\perp B = qvB\sin\theta$}
    {Magnetic force on a moving charge.}{\fgch{2}, \fgch{5}}
  \fgrow{\fgYes}{$F_\mathrm{c} = \dfrac{mv^2}{r}$}
    {Net force to the centre (gives $r$).}{\fgch{2}, \fgch{5}}
  \fgrow{\fgYes}{$K = \dfrac{1}{2}mv^2$}
    {Kinetic energy.}{\fgch{2}, \fgch{5}}
  \fgrow{\fgComb}{$v = \sqrt{\dfrac{2qV}{m}}$}
    {Speed after $V$, from rest.}{\fgch{2}}
\end{longtable}}

\section*{Same letter, different meaning}
{\small\renewcommand\arraystretch{1.1}\arrayrulecolor{hscrule}%
\noindent\begin{tabular}{@{}L{9mm}@{\hspace{1.5mm}}L{\dimexpr0.5\linewidth-12.5mm\relax}%
    @{\hspace{4mm}}L{9mm}@{\hspace{1.5mm}}L{\dimexpr0.5\linewidth-12.5mm\relax}@{}}
  $\lambda$ & wavelength, or the decay constant ($N_\mathrm{t} = N_0 e^{-\lambda t}$)
    & $T$ & temperature in K ($\lambda_\mathrm{max} = \frac{b}{T}$), or a period\\\hline
  $E$ & energy, or electric field ($E = \frac{V}{d}$, $v = \frac{E}{B}$)
    & $M$ & absolute magnitude; $M_\odot$: the Sun's mass\\\hline
  $e$ & elementary charge, or the number $2.718$ in $e^{-\lambda t}$
    & $m$ & mass, or the order of a line ($d\sin\theta = m\lambda$)\\\hline
  $A$ & mass number, or activity ($A = \lambda N$)
    & $n$ & level number, or a count (electrons, half-lives, gaps)\\\hline
  $N$ & number of neutrons, or number of nuclei
    & $R$ & Rydberg constant, or a radius\\\hline
  $Z$, $z$ & atomic number; redshift
    & u, $u$ & atomic mass unit; up quark; a speed\\\hline
  $V$, $v$ & voltage; speed
    & $\nu$, $v$ & neutrino (Greek nu); speed\\\hline
\end{tabular}\par}

\section*{Data sheet values used in this module}
The exam gives you the data sheet. These are the values that Module~\modulenumber\ uses,
exactly as printed.

{\small\renewcommand\arraystretch{1.1}\arrayrulecolor{hscrule}%
\noindent\begin{tabular}{@{}L{50mm}@{\hspace{3mm}}L{58mm}@{\hspace{3mm}}L{\dimexpr\linewidth-114mm\relax}@{}}
  \colhead{Quantity (as on the sheet)} & \colhead{Value} & \colhead{Used in}\\\hline
  Charge on electron & $q_\mathrm{e} = \qty{-1.602e-19}{C}$ & every chapter\\\hline
  Mass of electron & $m_\mathrm{e} = \qty{9.109e-31}{kg}$ & Chapters 1, 2, 3, 5\\\hline
  Mass of neutron & $m_\mathrm{n} = \qty{1.675e-27}{kg}$ & Chapters 2, 3\\\hline
  Mass of proton & $m_\mathrm{p} = \qty{1.673e-27}{kg}$ & Chapters 1, 2, 3, 5\\\hline
  Earth's gravitational acceleration & $g = \qty{9.8}{m.s^{-2}}$ & Chapter 2 (Millikan)\\\hline
  Speed of light & $c = \qty{3.00e8}{m.s^{-1}}$ & Chapters 1, 3, 4, 5\\\hline
  Electric permittivity constant & $\varepsilon_0 = \qty{8.854e-12}{A^2.s^4.kg^{-1}.m^{-3}}$
    & Chapter 2 (Coulomb)\\\hline
  Planck constant & $h = \qty{6.626e-34}{J.s}$ & Chapters 1, 3, 5\\\hline
  Rydberg constant & $R = \qty{1.097e7}{m^{-1}}$ & Chapter 3\\\hline
  Atomic mass unit & $u = \qty{1.661e-27}{kg} = \qty{931.5}{MeV}/c^2$ & Chapters 1, 2, 4\\\hline
  1 eV & $\qty{1.602e-19}{J}$ & every chapter\\\hline
  Wien's displacement constant & $b = \qty{2.898e-3}{m.K}$ & Chapter 1\\\hline
\end{tabular}\par}

\begin{itemize}
  \item The sheet prints the electron's charge as negative. Use its size,
        $e = \qty{1.602e-19}{C}$, and find directions separately.
  \item \emph{Not} on the data sheet (they are given when needed): $H_0$, the megaparsec,
        masses of nuclei in u, the density of oil, the laboratory wavelengths of spectral
        lines, the energy of one fission.
  \item Learn $E_n = -\dfrac{\qty{13.6}{eV}}{n^2}$ for hydrogen: it is not on the sheet
        (\cref{ch3:sec:bohr}).
  \item With masses in u, multiply the mass lost by 931.5 to get MeV. Do not mix this with
        kg and $c^2$ in one calculation: the two routes differ by about \qty{0.2}{\percent}.
  \item Use SI units in every formula, except the ``$\times\,931.5$'' rule: nm
        $\times\,10^{-9}$ gives m; \mbox{eV $\times\,\num{1.602e-19}$} gives J.
\end{itemize}
